\begin{frame}
    \frametitle{Once-through transitions provide expected demand of \gls{HALEU}}
        \begin{itemize}
            \item If we understand the demand for \gls{HALEU} for reactors, 
            then we can understand how much needs to be made and how that 
            demand affects the rest of the fuel cycle. 
            \pause
            \item Factors that affect demand:
            \begin{itemize}
                \item Reactor type
                \item Energy demand
            \end{itemize}
            \pause
            \item We can use fuel cycle simulators to model these transitions
        \end{itemize}
\end{frame}

\begin{frame}
    \frametitle{Fuel cycle models contains various assumptions}
    \begin{columns}
        
    \column[t]{6cm}
    \vspace{-0.9cm}
    \begin{figure}[t]
    \centering
    \begin{tikzpicture}[node distance=1cm]
        \node (mine) [agent] {Uranium Mine};
        \node (enrichment) [agent, below of=mine]{Enrichment};
        \node (reactor) [agent, below of=enrichment]{LWR};
        \node (adv_reactor) [transition, below of=enrichment, xshift=2cm]{Advanced Reactor};
        \node (wetstorage) [agent, below of=reactor]{Wet Storage};
        \node (drystorage) [agent, below of=wetstorage]{Dry Storage};
        \node (sinkhlw) [agent, below of=drystorage, xshift=0.75cm]{HLW Sink};
        \node (sinkllw) [agent, left of=enrichment,xshift=-1cm]{LLW Sink};
        \node (adv_rx_cooling) [transition, below of=adv_reactor]{Wet Storage};

        \draw [arrow] (mine) --  (enrichment); 
        \draw [arrow] (enrichment) -- (sinkllw);
        \draw [arrow] (enrichment) -- (reactor);
        \draw [arrow] (enrichment) -- (adv_reactor);
        \draw [arrow] (reactor) -- (wetstorage);
        \draw [arrow] (wetstorage) -- (drystorage);
        \draw [arrow] (drystorage) -- (sinkhlw);
        \draw [arrow] (adv_reactor) -- (adv_rx_cooling);
        \draw [arrow] (adv_rx_cooling) -- (sinkhlw);
        \end{tikzpicture}
    \caption{Fuel cycle facilities and material flow between facilities. Facilities 
    in red are deployed after 2025 for the transition. }
    \label{fig:fuel_cycle}
\end{figure}

        \column[t]{4.5cm}
        Model transitions using \Cyclus \cite{huff_fundamental_2016}
        \begin{itemize}
            \item Simulations model reactor deployment from 1965-2090
            \item<2-> \gls{LWR} commission dates are obtained from the IAEA PRIS
                database \cite{noauthor_power_1989}
            \item<2-> \glspl{LWR} are assumed to operate until their current license 
                expires \cite{noauthor_licensing_nodate}
            \item<3-> Transitions begin in 2025
            \item<3-> Manually calculate advanced reactor deployment
        \end{itemize}

\end{columns}
\end{frame}

\begin{frame}
    \frametitle{Advanced Reactors}
    \begingroup
        \renewcommand{\arraystretch}{1.5}
        \begin{table}
            \small
            \caption{Advanced reactor design specifications}
            \label{tab:reactor_summary}
            \begin{tabular}{ c c c c }
                \hline
                Design Criteria & USNC MMR 
                    \cite{noauthor_usnc_2021} & 
                    X-energy Xe-100 \cite{mulder_overview_2021} & 
                    VOYGR \cite{nuscale_chapter_2020-1,reyes_nuscale_2021,reyes_correction_2022}\\\hline
                Power Output (MWe) & 5 & 80 & 77\\
                Enrichment (\% $^{235}U$) & 19.75 & 15.5 & 4.09 \\
                Cycle Length (yr) & 20 & Online & 1.5 \\
                Number of cycles & 1 & 6 & 3\\
                Reactor Lifetime (yr) & 20 & 60 & 60\\
                Burnup ($\frac{MWd}{kg U}$) & 82 & 168 & 45\\
                \hline
            \end{tabular}
        \end{table}   
    \endgroup
\end{frame}

\begin{frame}
    \frametitle{Once-through scenario definitions}
        \begin{table}[ht]
            \centering
            \caption{Summary of the once-through fuel cycle transition scenarios.}
            \label{tab:scenarios_once-through}
            \begin{tabular}{c l l}
                    \hline
                    Scenario number & Reactors present & Energy growth\\\hline
                    1 & \glspl{LWR} & N/A \\
                    2 & \glspl{LWR} and \gls{MMR} & No growth \\
                    3 & \glspl{LWR} and Xe-100 & No growth \\
                    4 & \glspl{LWR}, Xe-100, and \gls{MMR}& No growth\\
                    5 & \glspl{LWR}, \gls{MMR}, and VOYGR & No growth\\
                    6 & \glspl{LWR}, Xe-100, and VOYGR & No growth\\
                    7 & \glspl{LWR}, Xe-100, \gls{MMR}, and VOYGR & No growth\\
                    8 & \glspl{LWR} and \gls{MMR}& 1\% growth \\
                    9 & \glspl{LWR} and Xe-100 & 1\% growth\\
                    10 & \glspl{LWR}, Xe-100, and \gls{MMR}& 1\% growth\\
                    11 & \glspl{LWR}, \gls{MMR}, and VOYGR & 1\% growth\\
                    12 & \glspl{LWR}, Xe-100, and VOYGR & 1\% growth\\
                    13 & \glspl{LWR}, Xe-100, \gls{MMR}, and VOYGR & 1\% growth\\
                    \hline
            \end{tabular}
        \end{table}
\end{frame}

\begin{frame}
    \frametitle{Advanced reactor deployment scheme}
    \begin{columns}
        \column[t]{5cm}
            Manually calculate the deployment scheme for advanced reactors
            \begin{itemize}
                \item Preferentially deploy reactors with larger power output first
                \item Deployment is based on the difference between the energy produced 
                      by reactors in the simulation and the demand at each time step
            \end{itemize}
        \column[t]{5cm}
            \begin{figure}
                \includegraphics[scale=0.25, trim=50 50 50 50,clip]{Deployment_Scheme.pdf}
                \caption{Example of how advanced reactors in Scenario 7 are deployed to 
                meet a demand of 530 MWe.}
                \label{fig:deployment}
            \end{figure}
                    
    \end{columns}
\end{frame}